\documentclass[a4paper,12pt]{article}
\usepackage[utf8]{inputenc} 
\usepackage{amsmath}
\usepackage[a4paper, margin=0.6in,top=0.3in, bottom=0.3in]{geometry} % Adjust the margin here
\usepackage{multirow}
\usepackage{nopageno}
\usepackage{graphicx} % Add this line to include graphics
\usepackage{tikz} % Add this line to use TikZ for drawing
\usepackage{tikz-3dplot} % Add this line to use 3D plotting
\setlength{\parindent}{0pt} % Set global no indent
    
\title{{\large Department of Physics, FNSPE, CTU in Prague} \\
                  02YELMA - Electricity and Magnetism \\ 
          {\Large \textbf{Midterm 2 (A - Solution)}}\vspace{-1cm}}
\date{}
\author{}

\begin{document}

\maketitle

\textbf{Q1:} By cylindrical symmetry, the magnetic field is purely azimuthal, $\vec B(s)=B_\phi(s)\,\hat\phi$. Using Ampère’s law for a circular loop of radius \(s\),

\[
B_\phi(s)(2\pi s)=\mu_0 I_{\text{enc}}
\]

We first relate \(J_0\) to the total current \(I\):
\[
I=\int_0^R J_0\left(\frac{s}{R}\right)(2\pi s\,ds) \Rightarrow \boxed{J_0=\frac{3I}{2\pi R^2}}
\]

For \(s<R\) (inside the conductor):

\[
I_{\text{enc}}(s)=\int_0^s J_0\left(\frac{s'}{R}\right)(2\pi s' \, ds')
=\frac{2\pi J_0 s^3}{3R}
\]
\[
B_\phi(s)(2\pi s)=\mu_0 \frac{2\pi J_0 s^3}{3R}
\]
\[
\boxed{
\vec B(s)=\frac{\mu_0 I}{2\pi R^3}s^2\,\hat\phi
\qquad (s<R)
}
\]

For \(s>R\) (outside the conductor):
\[
B_\phi(s)(2\pi s)=\mu_0 I
\]

\[
\boxed{
\vec B(s)=\frac{\mu_0 I}{2\pi s}\,\hat\phi
\qquad (s>R)
}
\]


\vspace*{0.3cm}
\textbf{Q2:} Between the plates, the changing electric field produces a magnetic field. Using Ampère-Maxwell law for a circular loop of radius \(s<R\),

\[
\oint \vec B\cdot d\vec l
=
\mu_0\epsilon_0
\frac{d\Phi_E}{dt}
\]
\[
B_\phi(s)(2\pi s)
=
\mu_0\epsilon_0
\frac{d}{dt}\left(E\pi s^2\right)
\]

\[
\vec B(s)
=
\frac{\mu_0\epsilon_0 s}{2}
\frac{dE}{dt}\,\hat\phi
\]
\[
\vec S
=
\frac{1}{\mu_0}\vec E\times\vec B =
-\frac{\epsilon_0 s}{2}
E\frac{dE}{dt}\,\hat s
\]
which points radially inward, indicating energy flows into the capacitor. Energy enters through the cylindrical side surface of area $dA=(2\pi R d)$

\[
P_{\text{in}}
=
\int \vec S\cdot d\vec a
=
\frac{\epsilon_0 R}{2}
E\frac{dE}{dt}(2\pi R d) = \epsilon_0\pi R^2 d\,
E\frac{dE}{dt} = \frac{dU}{dt}
\]

\[
\frac{dU}{dt}
=
\epsilon_0\pi R^2 d\,
E\frac{dE}{dt} \Rightarrow \boxed{
U(t)
=
\frac{1}{2}\epsilon_0\pi R^2 d\,E(t)^2
}
\]


\vspace*{0.3cm}
\textbf{Q3:} Since the wave propagates in free space, it travels in the \(+\hat z\) direction. The corresponding electric field is perpendicular to both \(\vec B\) and the direction of propagation:

\[
\vec E(z,t)=E_0\sin(kz-\omega t)\,\hat x, ~~E_0=cB_0,~~c=\frac{1}{\sqrt{\mu_0\epsilon_0}}
\]
(a) The instantaneous Poynting vector is

\[
\vec S
=
\frac{1}{\mu_0}\vec E\times\vec B, ~~ S=
\frac{1}{\mu_0}E_0B_0\sin^2(kz-\omega t) = \frac{cB_0^2}{\mu_0}\sin^2(kz-\omega t)
\]


\[
\boxed{
\langle S\rangle
=
\frac{cB_0^2}{2\mu_0}
}~~\text{since}~~\langle \sin^2(\cdot)\rangle=\frac{1}{2}.
\]

The area of the circular surface is $A=\pi a^2$. Therefore the average power carried by the wave is $P=\langle S\rangle A$
\[
\boxed{
P
=
\frac{cB_0^2}{2\mu_0}\,\pi a^2
}
\]

\bigskip

(b) The transmitted electric-field amplitude is $E_{\text{trans}}=E_0\cos\theta$ with $\theta=\frac{\pi}{4}$. Thus,

\[
E_{\text{trans}}
=
E_0\cos\left(\frac{\pi}{4}\right)
=
\frac{E_0}{\sqrt{2}}
\]

Since \(E_0=cB_0\), the transmitted magnetic-field amplitude is also reduced by the same factor:

\[
\boxed{
B_{\text{trans}}
=
\frac{B_0}{\sqrt{2}}
}
\]

\noindent

\vspace*{0.3cm}
\textbf{Q4:} Since the cylinder is metallic, the charge resides on its outer surface. The surface charge density is $\sigma=\frac{Q}{2\pi r h}$ because the curved surface area is $A=2\pi r h$ (neglecting the end caps). As the cylinder rotates, the surface charges move with tangential speed $v=\omega r$, thus the surface current density is

\[
K=\sigma v
=
\sigma \omega r = \frac{Q}{2\pi r h}\omega r = \frac{Q\omega}{2\pi h}
\]



If we treat the cylinder as stacked current loops, a thin strip of width \(dz\) carries current $dI=K\,dz$. Each strip forms a circular loop of area $A_{\text{loop}}=\pi r^2$ whose magnetic dipole moment is $d\mu=dI\cdot A_{\text{loop}}$. Therefore

\[
d\mu
=
K\,dz\cdot \pi r^2
\]

Integrating over the full height:

\[
\mu
=
\int_0^h K\pi r^2\,dz
=
K\pi r^2 h
\]

Substituting \(K\):

\[
\mu
=
\frac{Q\omega}{2\pi h}\cdot \pi r^2 h
\]

which gives

\[
\boxed{
\mu
=
\frac{1}{2}Q\omega r^2
} \Rightarrow \boxed{
\vec{\mu}
=
\frac{1}{2}Q\omega r^2\,\hat z
}
\]



\end{document}

